\documentclass[11pt,a4paper]{article}

% --- packages -----------------------------------------------------------------
\usepackage[utf8]{inputenc}
\usepackage[T1]{fontenc}
\usepackage{lmodern}
\usepackage[a4paper,margin=2.4cm,headheight=14pt]{geometry}
\usepackage{microtype}
\usepackage{xcolor}
\usepackage{listings}
\usepackage{hyperref}
\usepackage{enumitem}
\usepackage{tabularx}
\usepackage{booktabs}
\usepackage{longtable}
\usepackage{titlesec}
\usepackage{fancyhdr}
\usepackage{parskip}
\usepackage{amsmath,amssymb,amsthm}
\usepackage{tikz}
\usetikzlibrary{shapes,arrows.meta,positioning,calc,fit,decorations.pathreplacing}
\usepackage{tikz-cd}
\usepackage{algorithm}
\usepackage{algpseudocode}
\usepackage{caption}
\usepackage{subcaption}
\usepackage{multicol}
\usepackage{float}

% --- colours (Tokyo Night) ---------------------------------------------------
\definecolor{tnbg}{HTML}{1A1B26}
\definecolor{tnfg}{HTML}{C0CAF5}
\definecolor{tnblue}{HTML}{7AA2F7}
\definecolor{tnred}{HTML}{F7768E}
\definecolor{tnyellow}{HTML}{E0AF68}
\definecolor{tngreen}{HTML}{73DACA}
\definecolor{tncomment}{HTML}{565F89}
\definecolor{linkcolor}{HTML}{2C5AA0}
\definecolor{lightgray}{HTML}{EAEAEA}
\definecolor{nodecol}{HTML}{E8F0FE}
\definecolor{edgecol}{HTML}{1A73E8}
\definecolor{delcol}{HTML}{EA4335}
\definecolor{addcol}{HTML}{34A853}

% --- hyperref -----------------------------------------------------------------
\hypersetup{
  colorlinks=true,
  linkcolor=linkcolor,
  urlcolor=linkcolor,
  citecolor=linkcolor,
  pdftitle={sotBSD: A Capability-Secure Microkernel with Graph-Structured Filesystem},
  pdfauthor={Lucas Sotomayor (BASE4 Security)},
  pdfsubject={Microkernel OS, graph filesystem, formal verification},
  pdfkeywords={microkernel, capability, graph filesystem, DPO, treewidth, Ollivier-Ricci, deception, TLA+, Rust}
}

% --- listings -----------------------------------------------------------------
\lstdefinelanguage{rustlang}{
  keywords={fn,let,mut,pub,struct,enum,impl,trait,if,else,for,while,loop,match,return,use,mod,as,where,self,Self,const,static,unsafe,extern,move,async,await,box,break,continue,crate,dyn,false,true,in,ref,type},
  keywordstyle=\color{tnblue}\bfseries,
  ndkeywords={u8,u16,u32,u64,usize,i8,i16,i32,i64,isize,bool,char,str,Option,Result,Vec,Box,String},
  ndkeywordstyle=\color{tngreen}\bfseries,
  identifierstyle=\color{black},
  sensitive=true,
  comment=[l]{//},
  morecomment=[s]{/*}{*/},
  commentstyle=\color{tncomment}\itshape,
  stringstyle=\color{tnred},
  morestring=[b]",
  morestring=[b]',
}
\lstset{
  language=rustlang,
  basicstyle=\ttfamily\footnotesize,
  numbers=left,
  numberstyle=\tiny\color{tncomment},
  numbersep=8pt,
  frame=single,
  framesep=4pt,
  rulecolor=\color{lightgray},
  backgroundcolor=\color{lightgray!30},
  showstringspaces=false,
  breaklines=true,
  tabsize=2,
  captionpos=b,
}

% --- titlesec -----------------------------------------------------------------
\titleformat{\section}
  {\normalfont\Large\bfseries\color{tnblue}}
  {\thesection}{1em}{}
\titleformat{\subsection}
  {\normalfont\large\bfseries\color{tnblue!80!black}}
  {\thesubsection}{1em}{}
\titleformat{\subsubsection}
  {\normalfont\normalsize\bfseries}
  {\thesubsubsection}{1em}{}

% --- header / footer ----------------------------------------------------------
\pagestyle{fancy}
\fancyhf{}
\fancyhead[L]{\textit{sotBSD: Capability-Secure Microkernel with Graph-Structured FS}}
\fancyhead[R]{\textit{April 2026}}
\fancyfoot[C]{\thepage}
\renewcommand{\headrulewidth}{0.3pt}

% --- theorem environments -----------------------------------------------------
\theoremstyle{definition}
\newtheorem{definition}{Definition}[section]
\newtheorem{theorem}{Theorem}[section]
\newtheorem{corollary}[theorem]{Corollary}
\newtheorem{lemma}[theorem]{Lemma}
\newtheorem{proposition}[theorem]{Proposition}
\newtheorem{invariant}{Invariant}[section]
\newtheorem{constraint}{Constraint}[section]
\newtheorem{property}{Property}[section]

\theoremstyle{remark}
\newtheorem{remark}{Remark}[section]

% --- custom commands ----------------------------------------------------------
\newcommand{\TG}{\ensuremath{\mathcal{TG}}}
\newcommand{\sotfs}{\textsc{sotFS}}
\newcommand{\sotos}{\textsc{sotOS}}
\newcommand{\sotbsd}{\textsc{sotBSD}}
\newcommand{\DPO}{\textsc{dpo}}
\newcommand{\MSO}{\textsc{mso}}
\newcommand{\WAL}{\textsc{wal}}
\newcommand{\CDT}{\textsc{cdt}}
\newcommand{\NodeType}{\texttt{NodeType}}
\newcommand{\EdgeType}{\texttt{EdgeType}}
\newcommand{\tw}{\ensuremath{\mathrm{tw}}}
\newcommand{\OR}{\ensuremath{\kappa}}  % Ollivier-Ricci
\newcommand{\Wone}{\ensuremath{W_1}}
\newcommand{\cat}[1]{\ensuremath{\mathbf{#1}}}

% --- bibliography hack (manual) -----------------------------------------------
% We use manual \bibitem for maximum portability with pdflatex.

% --- title metadata -----------------------------------------------------------
\title{
  \vspace{-1.5cm}
  \textbf{sotBSD}: A Capability-Secure Microkernel with \\
  Graph-Structured Filesystem, Algebraic Verification, \\
  and Deception-in-Depth \\[0.6em]
  \normalsize Conference Submission Draft
}
\author{
  Lucas Sotomayor \\
  \texttt{lucas@base4.io} \\
  BASE4 Security
}
\date{April 2026}

% =============================================================================
\begin{document}
% =============================================================================

\maketitle
\thispagestyle{empty}

% =============================================================================
% ABSTRACT
% =============================================================================
\begin{abstract}
\noindent
We present \sotbsd{}, a capability-secure microkernel operating system
written entirely in Rust, featuring five kernel primitives and a
SOT (Secure Object Transaction) exokernel layer. Atop this foundation
we build \sotfs{}, the first filesystem whose metadata is modeled as a
typed graph and whose POSIX operations are defined as algebraic
Double Pushout (DPO) graph rewriting rules with formally verified
preconditions and gluing constraints. We prove that \sotfs{}'s type
graph maintains bounded treewidth---at most $\mathrm{LINK\_MAX} + O(1)$
in general, and at most 6 without hard links---enabling linear-time
evaluation of Monadic Second-Order (MSO) queries via Courcelle's
theorem. We introduce a novel application of Ollivier--Ricci discrete
curvature to filesystem graphs, enabling topology-aware anomaly
detection that identifies ransomware-like fan-out and exfiltration
staging with zero training data. Capability-gated graph projections
provide deception-in-depth: attacker-visible subgraphs are functorially
consistent yet strategically misleading. A four-layer verification
chain---TLA$^+$ model checking (24 invariants verified), crash
refinement proofs, Rust typestate enforcement, and a Coq extraction
roadmap---provides defense in depth against specification, protocol,
and implementation bugs. We evaluate \sotfs{} on 76 unit and integration
tests, 300{,}000 fuzz corpus runs with zero crashes, and 13 adversarial
attack scenarios with complete detection.
\end{abstract}

\tableofcontents
\newpage

% =============================================================================
\section{Introduction}
\label{sec:intro}
% =============================================================================

Modern operating systems manage filesystem metadata through implicit
graph structures---inode tables, directory entry lists, extent
trees---yet expose none of this structure for formal reasoning,
efficient querying, or security monitoring. The consequences are
well-documented: filesystems remain the largest source of
crash-consistency bugs~\cite{pillai2014}, access control is bolted on
via DAC or MAC rather than integrated into the data model, and anomaly
detection requires expensive per-file heuristics rather than
structural analysis.

This paper presents two tightly integrated contributions. First,
\sotbsd{}, a capability-secure microkernel with five primitives
(scheduling, IPC, capabilities, IRQ virtualization, frame allocation)
and a SOT exokernel providing transactional secure objects with
provenance tracking. Second, \sotfs{}, a graph-structured filesystem
built atop the SOT layer that represents metadata as a typed graph
$\TG = (V, E, \mathrm{src}, \mathrm{tgt}, \tau_V, \tau_E,
\mathrm{attr})$ with six node types and six edge types, and defines
all POSIX operations as Double Pushout (DPO) graph rewriting rules.

\paragraph{Thesis.}
By making filesystem metadata graph structure explicit and equipping
operations with algebraic semantics, we can simultaneously achieve
three goals that current systems treat as independent:
(1)~linear-time decidability of rich structural queries via bounded
treewidth and Courcelle's theorem,
(2)~topology-aware anomaly detection through Ollivier--Ricci discrete
curvature without training data, and
(3)~capability-gated deceptive projections that are functorially
consistent yet strategically misleading.

\paragraph{Contributions.}
This paper makes five numbered contributions:

\begin{enumerate}[label=\textbf{C\arabic*}.,leftmargin=2em]
\item \textbf{First graph-structured filesystem with algebraic (DPO)
  operational semantics.} Every POSIX metadata operation is a
  formally-specified graph rewriting rule with preconditions
  (application conditions), gluing constraints (dangling condition,
  identification condition), and postconditions expressed as typed
  graph morphisms. This provides a sound algebraic foundation for
  reasoning about operation composition, conflict detection, and
  crash recovery.

\item \textbf{Bounded treewidth preservation enabling linear-time MSO
  queries via Courcelle's theorem.} We prove that the \sotfs{} type
  graph has treewidth at most $\mathrm{LINK\_MAX} + O(1)$ in general,
  and at most~6 without hard links. Every DPO rule preserves this
  bound, enabling any MSO-definable property (reachability, acyclicity,
  permission consistency) to be checked in $O(n)$ time where $n$ is the
  number of graph nodes.

\item \textbf{Novel application of Ollivier--Ricci curvature to
  filesystem anomaly detection.} By computing discrete Ricci curvature
  on filesystem graph edges, we detect topological anomalies---sudden
  fan-out (ransomware), deep chain creation (directory traversal
  attacks), mass deletion---without training data, achieving 100\%
  detection on 13 adversarial scenarios.

\item \textbf{Capability-gated graph projections for
  deception-in-depth.} We define a deception functor
  $D_{\mathrm{FS}} : \cat{TGraph} \to \cat{TGraph}$ that produces
  capability-consistent yet strategically misleading subgraph views,
  providing active deception against post-exploitation reconnaissance.

\item \textbf{Four-layer verification chain:
  TLA$^+$ $\to$ crash refinement $\to$ Rust typestate $\to$ Coq
  (planned).} We verify 24 TLA$^+$ invariants across four
  specifications, implement Rust typestate machines that mirror TLA$^+$
  states, and outline the extraction path to Coq for end-to-end
  machine-checked proofs.
\end{enumerate}

The remainder of this paper is organized as follows.
Section~\ref{sec:background} reviews necessary background.
Section~\ref{sec:related} surveys related work.
Section~\ref{sec:arch} presents the \sotbsd{} kernel architecture.
Section~\ref{sec:sotfs} defines \sotfs{} formally.
Section~\ref{sec:deception} describes the deception engine.
Section~\ref{sec:verification} details the verification chain.
Section~\ref{sec:impl} covers implementation.
Section~\ref{sec:eval} presents evaluation results.
Section~\ref{sec:discussion} discusses limitations.
Section~\ref{sec:conclusion} concludes.

% =============================================================================
\section{Background}
\label{sec:background}
% =============================================================================

\subsection{Double Pushout Graph Rewriting}
\label{sec:bg-dpo}

Double Pushout (DPO) rewriting~\cite{ehrig2006} is the standard
algebraic framework for graph transformation. A \emph{production rule}
$p = (L \xleftarrow{l} K \xrightarrow{r} R)$ consists of three
graphs: a left-hand side $L$ (the pattern to match), a gluing graph
$K$ (the interface preserved during rewriting), and a right-hand
side $R$ (the replacement). Applying $p$ to a host graph $G$ via a
match $m : L \to G$ proceeds in two steps:

\begin{enumerate}
\item \textbf{Deletion.} Compute the pushout complement $D$ such that
  the left square of the following diagram is a pushout:
  \[
  \begin{tikzcd}
    L \arrow[d, "m"'] & K \arrow[l, "l"'] \arrow[r, "r"]
      \arrow[d, "d"] & R \arrow[d, "m^*"] \\
    G & D \arrow[l, "l^*"'] \arrow[r, "r^*"] & H
  \end{tikzcd}
  \]

\item \textbf{Addition.} Compute the pushout $H$ of
  $D \xleftarrow{d \circ l^{-1}} K \xrightarrow{r} R$, yielding the
  result graph~$H$.
\end{enumerate}

The pushout complement exists if and only if the \emph{gluing
condition} holds: (i)~the \emph{dangling condition}---no edge in
$G \setminus m(L \setminus l(K))$ is incident to a deleted
node---and (ii)~the \emph{identification condition}---$m$ does not
identify two elements of $L$ unless both are in~$l(K)$.

\subsection{Treewidth and Courcelle's Theorem}
\label{sec:bg-treewidth}

We recall the formal definition of tree decomposition and treewidth,
following Robertson and Seymour~\cite{robertson1986}.

\begin{definition}[Tree Decomposition]
\label{def:tree-decomp}
A \emph{tree decomposition} of a graph $G = (V, E)$ is a pair
$(T, \{X_t\}_{t \in V(T)})$ where $T$ is a tree and each node
$t \in V(T)$ is associated with a \emph{bag} $X_t \subseteq V$
satisfying:
\begin{enumerate}[label=(\roman*)]
\item \textbf{Vertex cover:} $\bigcup_{t \in V(T)} X_t = V$.
\item \textbf{Edge cover:} For every edge $(u, v) \in E$, there
  exists $t \in V(T)$ with $\{u, v\} \subseteq X_t$.
\item \textbf{Running intersection:} For every $v \in V$, the set
  $\{t \in V(T) \mid v \in X_t\}$ induces a connected subtree of $T$.
\end{enumerate}
The \emph{width} of the decomposition is $\max_{t \in V(T)} |X_t| - 1$.
\end{definition}

\begin{definition}[Treewidth]
\label{def:treewidth}
The \emph{treewidth} $\tw(G)$ of a graph $G$ is the minimum width
over all tree decompositions of $G$. Informally, it measures how
``tree-like'' a graph is: trees have treewidth~1, series-parallel
graphs have treewidth~2, and the complete graph $K_n$ has
treewidth~$n - 1$.
\end{definition}

\begin{theorem}[Courcelle~\cite{courcelle1990}]
\label{thm:courcelle}
For every MSO sentence $\varphi$ and every constant $k$, there exists
a linear-time algorithm that decides whether $\varphi$ holds on any
graph of treewidth at most~$k$.
\end{theorem}

The constant factor depends exponentially on $k$ and the quantifier
depth of $\varphi$, but for fixed formulas and bounded $k$ the
algorithm runs in $O(n)$ time. Bodlaender~\cite{bodlaender1998}
showed that a tree decomposition of width~$k$ can itself be computed
in $O(n)$ time for fixed~$k$.

\begin{definition}[MSO on Graphs]
\label{def:mso}
Monadic Second-Order logic on graphs extends first-order logic with
quantification over \emph{sets} of vertices and edges. An MSO formula
over a graph $G = (V, E)$ may use:
\begin{itemize}[leftmargin=1.4em]
\item Individual variables $x, y, \ldots$ ranging over $V \cup E$.
\item Set variables $X, Y, \ldots$ ranging over $2^V \cup 2^E$.
\item Predicates: $\mathrm{inc}(v, e)$ (vertex $v$ is incident to
  edge $e$), $v \in X$ (membership), equality.
\item Boolean connectives and quantifiers $\forall, \exists$ over
  both individual and set variables.
\end{itemize}
\end{definition}

MSO can express many filesystem properties of interest: ``Is the
directory subgraph acyclic?'', ``Can principal $P$ reach file $f$
through capabilities?'', ``Does every file in subtree $T$ satisfy
permission constraint $\phi$?''. By Theorem~\ref{thm:courcelle},
all such properties are decidable in $O(n)$ on \sotfs{} graphs.

\subsection{Ollivier--Ricci Curvature on Graphs}
\label{sec:bg-curvature}

Ollivier~\cite{ollivier2009} defined a notion of Ricci curvature for
metric spaces equipped with a family of probability measures. On a
graph $G = (V, E)$ with the shortest-path metric $d$, we assign to
each vertex $u$ a probability measure $\mu_u$ on its neighbors. The
\emph{Ollivier--Ricci curvature} of an edge $(u, v) \in E$ is:

\begin{equation}
\label{eq:ollivier-ricci}
\OR(u, v) = 1 - \frac{\Wone(\mu_u, \mu_v)}{d(u, v)}
\end{equation}

\noindent
where $\Wone(\mu_u, \mu_v)$ is the Wasserstein-1 (Earth Mover's)
distance between the measures $\mu_u$ and $\mu_v$.

Following Lin, Lu, and Yau~\cite{lin2011}, we use the \emph{lazy
random walk} with parameter $\alpha = 0.5$:
\[
\mu_u^{(\alpha)}(w) =
\begin{cases}
  \alpha & \text{if } w = u \\
  (1 - \alpha) / \deg(u) & \text{if } (u, w) \in E \\
  0 & \text{otherwise}
\end{cases}
\]

The Wasserstein-1 distance used in Equation~\eqref{eq:ollivier-ricci}
is defined via optimal transport:

\begin{definition}[Wasserstein-1 Distance]
\label{def:wasserstein}
Let $\mu, \nu$ be probability measures on a metric space $(X, d)$.
The Wasserstein-1 distance is:
\[
\Wone(\mu, \nu) = \inf_{\gamma \in \Gamma(\mu, \nu)}
  \int_{X \times X} d(x, y)\, d\gamma(x, y)
\]
where $\Gamma(\mu, \nu)$ is the set of all couplings (joint
distributions with marginals $\mu$ and $\nu$). On finite graphs,
this reduces to a linear program solvable in
$O(\deg(u) \cdot \deg(v))$ time for each edge $(u, v)$.
\end{definition}

Positive curvature $\OR > 0$ indicates clustering (neighbors of $u$
and $v$ overlap), while negative curvature $\OR < 0$ indicates
divergence (neighbors spread apart). In a filesystem graph, sudden
drops in curvature signal topological anomalies such as mass file
creation in a single directory (ransomware fan-out).

The key insight for filesystem monitoring is that curvature is a
\emph{local} quantity: $\OR(u, v)$ depends only on the 1-hop
neighborhoods of $u$ and $v$, making it efficient to compute
incrementally when the graph changes. Moreover, curvature captures
structural properties that degree alone cannot: a directory with 100
children in a flat hierarchy has different curvature than one with 100
children organized into 10 subdirectories, even though the root has
the same degree in both cases.

\subsection{Capability-Based Security}
\label{sec:bg-caps}

A \emph{capability} is an unforgeable reference to a kernel object
paired with an access rights mask~\cite{dennis1966}. The
\sotbsd{} capability model draws on seL4~\cite{klein2009},
Capsicum~\cite{watson2010}, and CHERI~\cite{woodruff2014}: capabilities
are the \emph{only} mechanism for accessing kernel objects, forming a
Capability Derivation Tree (\CDT{}) where rights can only be
attenuated (never amplified) through derivation.

\subsection{Microkernel Design Principles}
\label{sec:bg-ukernel}

The \sotbsd{} kernel follows the L4 tradition~\cite{liedtke1993}:
minimize the trusted computing base (TCB) to scheduling, IPC, and
address-space management, delegating all policy---including virtual
memory management, filesystems, and device drivers---to userspace
servers communicating via synchronous IPC.

% =============================================================================
\section{Related Work}
\label{sec:related}
% =============================================================================

\subsection{Verified Filesystems}

FSCQ~\cite{chen2015} pioneered machine-checked filesystem
verification in Coq, proving crash safety for a simple filesystem
with write-ahead logging. Yggdrasil~\cite{sigurbjarnarson2016}
automated crash-refinement proofs using SMT solvers, reducing
proof effort substantially. SquirrelFS~\cite{kadekodi2024} leveraged
Rust's type system to enforce crash consistency through typestate
patterns without external provers. Perennial~\cite{chajed2019}
introduced a concurrent separation logic framework for verifying
crash-safe distributed systems, including the GoJournal filesystem.

\sotfs{} differs from all of these in its data model: rather than
modeling filesystem state as flat maps or trees, we use typed graphs
with DPO rewriting rules. This provides a uniform algebraic framework
for both correctness proofs and runtime monitoring. Our verification
is currently less deep than FSCQ's machine-checked proofs but covers
a broader surface through the combination of TLA$^+$ model checking,
property-based testing, and fuzzing.

\subsection{Graph Databases and Graph-Structured Storage}

Graph databases such as Neo4j and systems like GraphChi process
graph-structured data efficiently but are designed for application-level
data, not filesystem metadata. The POSIX filesystem interface assumes
a hierarchical namespace, making it non-trivial to expose the
underlying graph structure without breaking compatibility.
BetrFS~\cite{jannen2015} uses write-optimized B$^\varepsilon$-trees
for metadata but does not expose graph structure.

\subsection{Filesystem Anomaly Detection}

PASS~\cite{muniswamy2006} pioneered filesystem provenance tracking,
recording causal chains of file operations. However, PASS captures
provenance as a separate artifact rather than integrating it into the
filesystem's data model. Subsequent work on intrusion
detection~\cite{king2003,goel2005} used system-call traces rather
than structural properties. Our approach is fundamentally different:
we compute anomaly scores from the \emph{topology} of the live
filesystem graph, requiring no training data and no separate
provenance store.

\subsection{Deception and Moving Target Defense}

HoneyD~\cite{provos2004} created network-level decoys.
MTD frameworks~\cite{jajodia2011} randomize attack surfaces.
Capability-based deception---where the security principal's
\emph{view} of the filesystem is a controlled projection of the true
state---has not been explored in the literature to our knowledge.
\sotfs{}'s deception functor provides the first formal framework
for capability-gated filesystem deception with functorial consistency
guarantees.

\subsection{Microkernel Operating Systems}

seL4~\cite{klein2009} established the gold standard for microkernel
verification. Genode~\cite{feske2015} provided a practical capability
framework. The L4 family~\cite{liedtke1993} demonstrated that
microkernel IPC overhead can be minimized. Rump kernels~\cite{kantee2012}
showed that NetBSD kernel components can run as userspace libraries.
\sotbsd{} integrates ideas from all of these: L4-style synchronous
IPC, seL4-inspired capabilities with a CDT, Genode-like component
isolation, and rump-kernel personality layers.

% =============================================================================
\section{Architecture}
\label{sec:arch}
% =============================================================================

\subsection{Kernel Primitives}
\label{sec:arch-primitives}

The \sotbsd{} microkernel provides exactly five primitives:

\begin{definition}[Kernel Primitive Set]
\label{def:primitives}
The \sotbsd{} kernel implements the primitive set
$\mathcal{P} = \{\textsc{Sched}, \textsc{Ipc}, \textsc{Cap},
\textsc{Irq}, \textsc{Frame}\}$ where:
\begin{itemize}[leftmargin=1.4em]
\item $\textsc{Sched}$: Preemptive 4-level priority MLQ scheduler
  with per-core run queues, Chase--Lev work stealing, and EDF
  deadline scheduling.
\item $\textsc{Ipc}$: L4-style synchronous endpoints (rendezvous)
  and asynchronous channels (ring buffers).
\item $\textsc{Cap}$: Capability table with generation counters
  (ABA-safe) and a Capability Derivation Tree (\CDT{}) supporting
  transitive revocation.
\item $\textsc{Irq}$: IRQ virtualization mapping hardware interrupts
  to IPC notifications.
\item $\textsc{Frame}$: Physical frame allocator (bitmap-based,
  $O(1)$ amortized) with per-process resource limits.
\end{itemize}
\end{definition}

All objects are accessed exclusively through capabilities. The kernel
maintains no global namespace, no process table visible to userspace,
and no filesystem state. Virtual memory management is delegated to a
userspace VMM server.

\begin{property}[TCB Minimality]
\label{prop:tcb}
The kernel TCB consists of approximately 10{,}000 lines of Rust.
Only code executing in Ring~0 can access capability tables, modify
page tables, or configure interrupt routing. All other functionality
---including virtual memory management, filesystem operations, device
drivers, and network stacks---executes in Ring~3 userspace processes.
\end{property}

The scheduler implements a four-level Multi-Level Queue (MLQ) with
priorities: \textsc{Realtime} (level~0), \textsc{High} (level~1),
\textsc{Normal} (level~2), and \textsc{Low} (level~3). Each CPU
maintains a per-core run queue, and idle CPUs perform Chase--Lev
work stealing from the busiest queue. Resource limits (CPU ticks
and memory pages) are enforced per-process in both the scheduler
and the frame allocator, preventing resource exhaustion attacks.

\begin{property}[Scheduling Fairness]
\label{prop:sched-fairness}
Within each priority level, the scheduler guarantees that every
ready thread receives CPU time within a bounded number of scheduling
quanta. Across priority levels, higher-priority threads preempt
lower-priority threads, with work stealing ensuring load balance
across cores.
\end{property}

\subsection{SOT Exokernel Layer}
\label{sec:arch-sot}

Above the five primitives, the SOT (Secure Object Transaction) layer
provides 12 syscalls (numbers 300--311) for transactional operations on
secure objects:

\begin{definition}[Secure Object]
\label{def:sot}
A \emph{secure object} $o = (\mathrm{id}, \tau, \mathrm{data},
\mathrm{prov})$ consists of a unique identifier, a type tag $\tau$, a
data payload, and a provenance record. Operations on secure objects are
mediated by capabilities and logged to a per-CPU provenance ring
(4096 entries, lock-free).
\end{definition}

The SOT layer supports single-tier commit, two-phase commit across
multiple objects, capability interposition with proxy domains, and
rollback to named snapshots. These transactions provide the foundation
for \sotfs{}'s crash-consistency guarantees.

\subsection{Address Space Isolation}
\label{sec:arch-as}

Each process executes in its own address space (separate CR3). The
kernel provides \texttt{SYS\_AS\_CREATE} and \texttt{SYS\_AS\_CLONE}
(with copy-on-write) syscalls. Page faults in user address spaces are
delivered to a registered VMM server via IPC, which resolves them by
mapping physical frames.

\begin{property}[Address Space Isolation]
\label{prop:as-isolation}
For any two processes $P_i, P_j$ with $i \neq j$, the set of physical
frames accessible to $P_i$ and $P_j$ are disjoint unless explicitly
shared through a capability-mediated mapping operation.
\end{property}

\subsection{IPC Architecture}
\label{sec:arch-ipc}

IPC uses two mechanisms:

\begin{enumerate}
\item \textbf{Synchronous endpoints} (L4-style rendezvous): sender
  blocks until a receiver is ready, then registers are transferred
  directly (zero-copy for register-sized messages, up to 8 registers
  $\times$ 8 bytes = 64 bytes per transfer). Call semantics
  (send + wait-for-reply) support RPC patterns.

\item \textbf{Asynchronous channels}: bounded ring buffers in shared
  memory for bulk data transfer. The kernel provides only the
  notification mechanism; the ring buffer protocol is implemented in
  userspace.
\end{enumerate}

A kernel-side service registry (syscalls 130--131) maps service names
to endpoint capabilities, enabling dynamic service discovery.

\begin{property}[IPC Delivery Guarantees]
\label{prop:ipc}
The synchronous endpoint protocol guarantees:
\begin{enumerate}[label=(\roman*)]
\item \textbf{No duplication:} Each message is delivered exactly once.
\item \textbf{No loss:} If both sender and receiver are live and the
  endpoint is not destroyed, delivery eventually occurs (under fair
  scheduling).
\item \textbf{Ordering:} Messages from a single sender to a single
  endpoint are delivered in FIFO order.
\item \textbf{Timeout support:} \texttt{SYS\_CALL\_TIMEOUT} (syscall
  135) allows callers to specify a deadline in scheduler ticks,
  preventing indefinite blocking on unresponsive servers.
\end{enumerate}
\end{property}

The IPC fast path transfers 8 registers (64 bytes) without any memory
copy, achieving sub-microsecond latency for small messages. Larger
payloads use shared-memory channels with kernel-mediated notification.
Capability transfer during IPC follows the grant/delegate model:
the sender can include a capability in the message, which the kernel
validates and installs in the receiver's capability table.

\subsection{Capability Model}
\label{sec:arch-caps}

\begin{definition}[Capability]
\label{def:cap}
A capability $c = (\mathrm{obj}, \mathrm{rights}, \mathrm{gen},
\mathrm{parent})$ references a kernel object $\mathrm{obj}$ with
rights mask $\mathrm{rights} \subseteq \mathcal{R}$, generation
counter $\mathrm{gen}$ for ABA safety, and parent index
$\mathrm{parent}$ in the \CDT{}.
\end{definition}

\begin{theorem}[Monotonic Attenuation]
\label{thm:monotonic}
For any capability derivation $c' = \textsc{Grant}(c, \mathrm{mask})$,
$\mathrm{rights}(c') = \mathrm{rights}(c) \cap \mathrm{mask}
\subseteq \mathrm{rights}(c)$.
\end{theorem}

\begin{proof}
By construction: the \textsc{Grant} operation computes the bitwise
AND of the parent's rights and the mask. Since $A \cap B \subseteq A$
for any sets $A, B$, the derived capability's rights are a subset of
the parent's. The \CDT{} records the parent--child relationship,
ensuring transitive revocation respects the derivation order.
\end{proof}

\begin{theorem}[Transitive Revocation]
\label{thm:revocation}
If $\textsc{Revoke}(c)$ is invoked, then $c$ and all capabilities
$c'$ such that $c \rightsquigarrow^* c'$ in the \CDT{} are marked
dead, and all subsequent access attempts via any $c'$ fail.
\end{theorem}

\begin{proof}
The \CDT{} is a forest (each capability has at most one parent).
Revocation performs a depth-first traversal of the subtree rooted at
$c$, marking each node dead. Since every capability check verifies
the \texttt{alive} flag before granting access, and the flag
transitions are monotonic (alive $\to$ dead, never reversed), no
revoked capability can be used after revocation completes.
\end{proof}

% =============================================================================
\section{\sotfs{}: Graph-Structured Filesystem}
\label{sec:sotfs}
% =============================================================================

\subsection{Type Graph}
\label{sec:sotfs-typegraph}

\begin{definition}[Type Graph]
\label{def:typegraph}
The \sotfs{} \emph{type graph} is a tuple
$\TG = (V, E, \mathrm{src}, \mathrm{tgt}, \tau_V, \tau_E,
\mathrm{attr})$ where:
\begin{itemize}[leftmargin=1.4em]
\item $V$ is the set of vertices (metadata nodes).
\item $E \subseteq V \times V$ is the set of directed edges.
\item $\mathrm{src}, \mathrm{tgt} : E \to V$ are source and target
  functions.
\item $\tau_V : V \to \NodeType$ assigns each vertex a type from
  $\NodeType = \{\texttt{Inode}, \texttt{Directory},
  \texttt{Capability}, \texttt{Transaction}, \texttt{Version},
  \texttt{Block}\}$.
\item $\tau_E : E \to \EdgeType$ assigns each edge a type from
  $\EdgeType = \{\texttt{contains}, \texttt{grants},
  \texttt{delegates}, \texttt{derivesFrom}, \texttt{supersedes},
  \texttt{pointsTo}\}$.
\item $\mathrm{attr} : V \cup E \to \mathcal{A}$ assigns
  attributes (sizes, timestamps, permissions, reference counts)
  to nodes and edges.
\end{itemize}
\end{definition}

The type constraints on edges enforce structural invariants:

\begin{constraint}[Edge Typing]
\label{con:edge-typing}
Each edge type has fixed source and target node types:
\begin{center}
\begin{tabular}{@{}lll@{}}
\toprule
\textbf{Edge Type} & \textbf{Source Type} & \textbf{Target Type} \\
\midrule
\texttt{contains}    & \texttt{Directory}   & \texttt{Inode} $\cup$ \texttt{Directory} \\
\texttt{grants}      & \texttt{Capability}  & \texttt{Inode} $\cup$ \texttt{Directory} \\
\texttt{delegates}   & \texttt{Capability}  & \texttt{Capability} \\
\texttt{derivesFrom} & \texttt{Version}     & \texttt{Version} \\
\texttt{supersedes}  & \texttt{Transaction} & \texttt{Transaction} \\
\texttt{pointsTo}    & \texttt{Inode}       & \texttt{Block} \\
\bottomrule
\end{tabular}
\end{center}
\end{constraint}

\subsection{DPO Rewriting Rules}
\label{sec:sotfs-dpo}

Every POSIX metadata operation in \sotfs{} is specified as a DPO
graph rewriting rule. We present three representative rules.

\subsubsection{Rule: \texttt{create\_file}}

The \texttt{create\_file} rule creates a new inode $i$ and a
\texttt{contains} edge from a parent directory $d$ to $i$:

\begin{equation}
\label{eq:create-file}
\begin{tikzcd}[column sep=2.5em]
L \arrow[d, "m"'] & K \arrow[l, "l"'] \arrow[r, "r"]
  \arrow[d, "d"] & R \arrow[d, "m^*"] \\
G & D \arrow[l, "l^*"'] \arrow[r, "r^*"] & H
\end{tikzcd}
\end{equation}

\noindent
where:
\begin{align*}
L &= \{d : \texttt{Dir}\} &
K &= \{d : \texttt{Dir}\} \\
R &= \{d : \texttt{Dir},\; i : \texttt{Inode},\;
       (d \xrightarrow{\texttt{contains}} i)\}
\end{align*}

The match $m$ identifies $d$ in the host graph $G$. The gluing graph
$K$ preserves $d$, and $R$ adds the new inode $i$ and the
\texttt{contains} edge. The application condition requires that no
existing child of $d$ has the same name (unique name invariant).

\subsubsection{Rule: \texttt{unlink}}

The \texttt{unlink} rule removes a \texttt{contains} edge from
directory $d$ to target $t$:

\begin{align*}
L &= \{d : \texttt{Dir},\; t : \texttt{Inode},\;
       (d \xrightarrow{\texttt{contains}} t)\} \\
K &= \{d : \texttt{Dir},\; t : \texttt{Inode}\} \\
R &= \{d : \texttt{Dir},\; t : \texttt{Inode}\}
\end{align*}

The dangling condition is critical here: if $t$ has
$\mathrm{link\_count} = 1$, then $t$'s \texttt{pointsTo} edges to
blocks must also be removed (cascading deletion). If
$\mathrm{link\_count} > 1$, only the directory entry is removed and
$t$ persists.

\subsubsection{Rule: \texttt{rename}}

The \texttt{rename} rule atomically removes a \texttt{contains} edge
from source directory $s$ to target $t$ and adds a \texttt{contains}
edge from destination directory $d'$ to $t$:

\begin{align*}
L &= \{s : \texttt{Dir},\; t : \texttt{Inode/Dir},\;
       (s \xrightarrow{\texttt{contains}} t)\} \\
K &= \{t : \texttt{Inode/Dir}\} \\
R &= \{d' : \texttt{Dir},\; t : \texttt{Inode/Dir},\;
       (d' \xrightarrow{\texttt{contains}} t)\}
\end{align*}

An additional application condition prevents creating directory cycles
when $t$ is a directory and $d'$ is a descendant of $t$.

\subsubsection{Rule: \texttt{hard\_link}}

The \texttt{hard\_link} rule adds a second \texttt{contains} edge
from a directory $d'$ to an existing inode $i$ that is already
reachable through another directory $d$:

\begin{align*}
L &= \{d' : \texttt{Dir},\; i : \texttt{Inode},\;
       d : \texttt{Dir},\;
       (d \xrightarrow{\texttt{contains}} i)\} \\
K &= \{d' : \texttt{Dir},\; i : \texttt{Inode},\;
       d : \texttt{Dir},\;
       (d \xrightarrow{\texttt{contains}} i)\} \\
R &= \{d' : \texttt{Dir},\; i : \texttt{Inode},\;
       d : \texttt{Dir},\;
       (d \xrightarrow{\texttt{contains}} i),\;
       (d' \xrightarrow{\texttt{contains}} i)\}
\end{align*}

Application conditions: (i)~$\tau_V(i) \neq \texttt{Directory}$
(POSIX prohibits hard links to directories), (ii)~the name for the
new link is unique within $d'$, and (iii)~$\mathrm{attr}(i).\mathrm{nlink}
< \mathrm{LINK\_MAX}$. Note that this rule increases the treewidth
of the graph because $i$ now appears in the bags of both $d$ and
$d'$ in any tree decomposition.

\subsubsection{Rule: \texttt{mkdir}}

The \texttt{mkdir} rule creates a new directory node $d'$ and
a \texttt{contains} edge from parent $d$ to $d'$:

\begin{align*}
L &= \{d : \texttt{Dir}\} \\
K &= \{d : \texttt{Dir}\} \\
R &= \{d : \texttt{Dir},\; d' : \texttt{Dir},\;
       (d \xrightarrow{\texttt{contains}} d')\}
\end{align*}

The rule mirrors \texttt{create\_file} but creates a
\texttt{Directory} node instead of an \texttt{Inode}. The new
directory is initialized with self-referential ``.'' and parent
``..'' entries (modeled as attributes rather than explicit edges
to avoid cycles in the type graph).

\subsubsection{Composition and Confluence}

A key property of the DPO framework is the ability to reason about
rule composition. Two rules $p_1$ and $p_2$ are \emph{parallel
independent} if neither deletes a graph element required by the
other's match. For \sotfs{}, most rules operating on disjoint
subtrees are parallel independent, enabling concurrent application
in a multi-threaded environment.

\begin{lemma}[Parallel Independence]
\label{lem:parallel-indep}
Two \sotfs{} DPO rule applications $p_1$ at match $m_1$ and $p_2$
at match $m_2$ are parallel independent if
$m_1(L_1 \setminus l_1(K_1)) \cap m_2(L_2) = \emptyset$ and
$m_2(L_2 \setminus l_2(K_2)) \cap m_1(L_1) = \emptyset$. In this
case, both orderings produce the same result graph $H$.
\end{lemma}

\begin{proof}
Standard result from algebraic graph
transformation~\cite{ehrig2006}, Theorem 5.21 (Parallelism Theorem).
The conditions ensure that neither rule deletes or modifies elements
matched by the other.
\end{proof}

\subsection{Treewidth Bound}
\label{sec:sotfs-treewidth}

We now prove the central structural theorem of \sotfs{}.

\begin{theorem}[Treewidth Bound]
\label{thm:treewidth}
For any \sotfs{} type graph instance $G$ conforming to the type
constraints of Definition~\ref{def:typegraph}:
\[
\tw(G) \leq \mathrm{LINK\_MAX} + O(1)
\]
Without hard links (i.e., every inode has at most one incoming
\texttt{contains} edge):
\[
\tw(G) \leq 6
\]
\end{theorem}

\begin{proof}
We construct an explicit tree decomposition.

\textbf{Case 1: No hard links.} The directory hierarchy forms a
forest (rooted tree when single-rooted). Each non-leaf vertex in the
tree decomposition bag contains at most:
\begin{itemize}[leftmargin=1.4em]
\item 1 directory node,
\item 1 parent directory node (for the \texttt{contains} edge),
\item up to 1 capability node (for the \texttt{grants} edge),
\item up to 1 transaction node (for transactional wrapping),
\item up to 1 version node (for versioning),
\item up to 1 block node (for the \texttt{pointsTo} edge from an
  inode to its first block).
\end{itemize}
Thus each bag has at most 7 vertices (including the node itself),
giving width $\leq 6$.

\textbf{Case 2: With hard links.} An inode $i$ with $h$ hard links
has $h$ incoming \texttt{contains} edges from (possibly distinct)
directories $d_1, \ldots, d_h$. In the tree decomposition, $i$ must
appear in the bags of all $d_j$, creating a ``hub'' of width $h + O(1)$.
POSIX limits $h \leq \mathrm{LINK\_MAX}$ (typically 65{,}535 on ext4),
so $\tw(G) \leq \mathrm{LINK\_MAX} + O(1)$.

\textbf{Preservation under DPO rules.} Each DPO rule either
(a)~adds a leaf node with bounded degree (create, mkdir, hard\_link
with $h < \mathrm{LINK\_MAX}$), (b)~removes edges without increasing
degree (unlink, rmdir), or (c)~moves an edge between existing nodes
(rename) without increasing any node's degree. Therefore, if
$\tw(G) \leq k$ before the rule application, then $\tw(H) \leq k$
after.
\end{proof}

\begin{corollary}[Linear-Time MSO Queries]
\label{cor:mso}
For any fixed MSO sentence $\varphi$ and any \sotfs{} instance $G$
without hard links, $\varphi$ can be evaluated on $G$ in time $O(|V|)$.
\end{corollary}

\begin{proof}
Immediate from Theorem~\ref{thm:treewidth} ($\tw \leq 6$) and
Courcelle's theorem (Theorem~\ref{thm:courcelle}).
\end{proof}

\begin{remark}
Practical MSO queries on \sotfs{} include: ``Is the directory
hierarchy acyclic?'' (safety), ``Can process $P$ reach file $f$
through any chain of capabilities?'' (authorization), and ``Do all
files in directory $d$ have consistent permissions?'' (policy
compliance). All are decidable in $O(n)$ on \sotfs{} graphs.
\end{remark}

\subsection{Structural Invariants}
\label{sec:sotfs-invariants}

\sotfs{} maintains the following invariants, each verified by TLA$^+$
model checking:

\begin{invariant}[Unique Names]
\label{inv:unique-names}
For any directory node $d$, no two outgoing \texttt{contains} edges
carry the same name attribute: $\forall d, \forall e_1, e_2 \in
\mathrm{out}(d, \texttt{contains})$, $e_1 \neq e_2 \implies
\mathrm{attr}(e_1).\mathrm{name} \neq \mathrm{attr}(e_2).\mathrm{name}$.
\end{invariant}

\begin{invariant}[No Dangling Edges]
\label{inv:no-dangling}
Every edge endpoint exists in $V$: $\forall e \in E$,
$\mathrm{src}(e) \in V \land \mathrm{tgt}(e) \in V$.
\end{invariant}

\begin{invariant}[Link Count Consistency]
\label{inv:link-count}
For every inode $i$, $\mathrm{attr}(i).\mathrm{nlink} =
|\{e \in E \mid \tau_E(e) = \texttt{contains} \land
\mathrm{tgt}(e) = i\}|$.
\end{invariant}

\begin{invariant}[Directory Acyclicity]
\label{inv:acyclicity}
The subgraph induced by \texttt{contains} edges between
\texttt{Directory} nodes is a forest (directed acyclic graph with
unique paths from roots).
\end{invariant}

\begin{invariant}[Block Reference Count]
\label{inv:block-refcount}
For every block node $b$, the reference count equals the number of
incoming \texttt{pointsTo} edges: $\mathrm{attr}(b).\mathrm{refcount}
= |\{e \in E \mid \tau_E(e) = \texttt{pointsTo} \land
\mathrm{tgt}(e) = b\}|$.
\end{invariant}

\subsection{Ollivier--Ricci Curvature Monitor}
\label{sec:sotfs-curvature}

We define a continuous monitor on the \sotfs{} graph that computes
Ollivier--Ricci curvature to detect topological anomalies.

\begin{definition}[Filesystem Curvature]
\label{def:fs-curvature}
For each edge $(u, v) \in E$ of the \sotfs{} graph, the
\emph{filesystem curvature} is:
\[
\OR_{\mathrm{FS}}(u, v) = 1 - \frac{\Wone(\mu_u^{(0.5)}, \mu_v^{(0.5)})}{d(u, v)}
\]
where $\mu_u^{(0.5)}$ is the lazy random walk with $\alpha = 0.5$
(Section~\ref{sec:bg-curvature}) restricted to the undirected version
of the \sotfs{} graph, and $d$ is the shortest-path metric.
\end{definition}

\begin{theorem}[Anomaly Detection Soundness]
\label{thm:anomaly}
Let $G_t$ denote the \sotfs{} graph at time $t$, and let
$\Delta\OR(e) = \OR_t(e) - \OR_{t-1}(e)$ for edge $e$.
If a batch of $n$ files is created in a single directory $d$ at
time $t$, then:
\[
\Delta\OR(d, \mathrm{parent}(d)) \leq -\Omega\!\left(\frac{n}{\deg(d) + n}\right)
\]
That is, fan-out creation produces a measurable curvature drop
proportional to the ratio of new files to existing degree.
\end{theorem}

\begin{proof}[Proof sketch]
Before the batch, $d$ has degree $k$. After adding $n$ children,
$\deg(d) = k + n$. The lazy random walk from $d$ now spreads its
$(1-\alpha)$ mass over $k + n$ neighbors instead of $k$. The
Wasserstein distance $\Wone(\mu_d, \mu_{\mathrm{parent}(d)})$
increases because the mass from $d$ is more diffuse, while the
parent's measure is relatively unchanged. A direct computation
using the dual formulation of $\Wone$ yields the stated bound.
\end{proof}

\begin{property}[Ransomware Detection]
\label{prop:ransomware}
A ransomware-like attack creating $n \geq 100$ encrypted copies in a
single directory produces $\Delta\OR < -0.3$, which exceeds the
alert threshold $\theta = -0.15$ by a factor of~2, enabling detection
within a single monitoring epoch.
\end{property}

% =============================================================================
\section{Deception Engine}
\label{sec:deception}
% =============================================================================

\subsection{Deception Functor}
\label{sec:deception-functor}

We formalize deception as a functor on the category of typed graphs.

\begin{definition}[Category $\cat{TGraph}$]
\label{def:cat-tgraph}
The category $\cat{TGraph}$ has typed graphs
(Definition~\ref{def:typegraph}) as objects and typed graph morphisms
(node-type-preserving and edge-type-preserving graph homomorphisms) as
arrows.
\end{definition}

\begin{definition}[Deception Functor]
\label{def:deception-functor}
A \emph{deception functor} $D_{\mathrm{FS}} : \cat{TGraph} \to
\cat{TGraph}$ maps each typed graph $G$ to a ``deceptive projection''
$D_{\mathrm{FS}}(G)$ satisfying:
\begin{enumerate}[label=(\roman*)]
\item \textbf{Capability consistency:} For every capability $c$ in
  $D_{\mathrm{FS}}(G)$, the access rights and derivation relationships
  are consistent with those in $G$. That is, $D_{\mathrm{FS}}$
  preserves the \CDT{} structure restricted to visible capabilities.

\item \textbf{Type preservation:} $D_{\mathrm{FS}}$ preserves node
  types and edge types: if $(u, v)$ is a \texttt{contains} edge in
  $D_{\mathrm{FS}}(G)$, then $u$ is a \texttt{Directory} and $v$ is an
  \texttt{Inode} or \texttt{Directory}.

\item \textbf{Strategic misleading:} $D_{\mathrm{FS}}(G)$ may add
  \emph{honeypot} nodes (fake files, directories, credentials) and
  remove or rename real nodes, subject to (i) and (ii).
\end{enumerate}
\end{definition}

\begin{theorem}[Functorial Consistency]
\label{thm:functorial}
$D_{\mathrm{FS}}$ preserves identity morphisms and composition:
$D_{\mathrm{FS}}(\mathrm{id}_G) = \mathrm{id}_{D_{\mathrm{FS}}(G)}$
and $D_{\mathrm{FS}}(g \circ f) = D_{\mathrm{FS}}(g) \circ
D_{\mathrm{FS}}(f)$ for composable morphisms $f, g$.
\end{theorem}

\begin{proof}
Identity preservation holds because the deception policy is
state-dependent: the identity morphism (no change to $G$) produces
no change to $D_{\mathrm{FS}}(G)$. Composition preservation holds
because $D_{\mathrm{FS}}$ is defined pointwise on nodes and edges:
the image of a composed morphism is the composition of images.
\end{proof}

\subsection{Capability-Gated Projections}
\label{sec:deception-projections}

Each security domain $\sigma$ is associated with a capability set
$C_\sigma$ determining which subgraph of $\TG$ is visible. The
deception engine intercedes at the capability level:

\begin{definition}[Deceptive View]
\label{def:deceptive-view}
For domain $\sigma$ with capability set $C_\sigma$ and deception
policy $\pi_\sigma$, the \emph{deceptive view} is:
\[
V_\sigma = \mathrm{Reachable}(\TG, C_\sigma) \cup
  \mathrm{Honeypots}(\pi_\sigma) \setminus
  \mathrm{Hidden}(\pi_\sigma)
\]
where $\mathrm{Reachable}(\TG, C_\sigma)$ is the subgraph reachable
through capabilities in $C_\sigma$, $\mathrm{Honeypots}(\pi_\sigma)$
is the set of decoy nodes injected by policy $\pi_\sigma$, and
$\mathrm{Hidden}(\pi_\sigma)$ is the set of real nodes concealed
by $\pi_\sigma$.
\end{definition}

\subsection{Deception Profiles}
\label{sec:deception-profiles}

\sotbsd{} ships with four deception profiles:

\begin{enumerate}
\item \textbf{Ubuntu Web Server:} Projects a fake
  \texttt{/var/www/html} with honeypot PHP files, a decoy
  \texttt{/etc/shadow}, and a fake \texttt{.ssh/authorized\_keys}.

\item \textbf{CentOS Database:} Projects fake MySQL data directories,
  decoy \texttt{/etc/my.cnf}, and honeypot backup files.

\item \textbf{IoT Camera:} Minimal filesystem with fake firmware
  update endpoints and decoy RTSP configuration.

\item \textbf{Windows File Server:} Projects SMB shares with honeypot
  documents and fake Active Directory configuration files.
\end{enumerate}

\subsection{Categorical Framework}
\label{sec:deception-categories}

The full deception framework involves six categories and their
connecting functors:

\begin{definition}[Category System]
\label{def:cat-system}
The \sotbsd{} deception framework comprises:
\begin{enumerate}[label=(\roman*)]
\item $\cat{SO}$: Secure Objects. Objects are SOT entries; arrows are
  capability-mediated access operations.
\item $\cat{Dom}$: Security Domains. Objects are process domains;
  arrows are domain transitions (e.g., via \texttt{execve}).
\item $\cat{Chan}$: IPC Channels. Objects are endpoints/channels;
  arrows are message transmissions.
\item $\cat{Tx}$: Transactions. Objects are transaction states; arrows
  are commit/rollback operations.
\item $\cat{Prov}$: Provenance. Objects are provenance records; arrows
  are causal chains.
\item $\cat{TGraph}$: Typed Graphs (Definition~\ref{def:cat-tgraph}).
\end{enumerate}
\end{definition}

Key functors connecting these categories:

\begin{itemize}[leftmargin=1.4em]
\item $\mathrm{Access} : \cat{SO} \to \cat{Dom}$ maps each secure
  object access to the domain performing it.
\item $\mathrm{Provenance} : \cat{SO} \to \cat{Prov}$ is a
  \emph{faithful} functor recording the causal history of every object
  operation.
\item $\mathrm{Graph} : \cat{SO} \to \cat{TGraph}$ projects the
  current secure object state into the type graph representation.
\item $D : \cat{SO} \to \cat{SO}$ is the \emph{object-level deception
  endofunctor} that remaps object identifiers and injects decoys.
\item $D_{\mathrm{FS}} : \cat{TGraph} \to \cat{TGraph}$ is the
  graph-level deception functor (Definition~\ref{def:deception-functor}).
\end{itemize}

\begin{proposition}[Deception Commutation]
\label{prop:deception-commutes}
The following diagram commutes:
\[
\begin{tikzcd}
\cat{SO} \arrow[r, "D"] \arrow[d, "\mathrm{Graph}"'] &
  \cat{SO} \arrow[d, "\mathrm{Graph}"] \\
\cat{TGraph} \arrow[r, "D_{\mathrm{FS}}"'] & \cat{TGraph}
\end{tikzcd}
\]
That is, $\mathrm{Graph} \circ D = D_{\mathrm{FS}} \circ \mathrm{Graph}$.
\end{proposition}

% =============================================================================
\section{Formal Verification}
\label{sec:verification}
% =============================================================================

\subsection{Layer 1: TLA$^+$ Model Checking}
\label{sec:verif-tla}

We maintain four TLA$^+$ specifications covering the \sotfs{} graph
invariants, transaction protocol, capability system, and crash
recovery. Table~\ref{tab:tla} summarizes the verified properties.

\begin{table}[H]
\centering
\caption{TLA$^+$ specifications and verified invariants.}
\label{tab:tla}
\small
\begin{tabular}{@{}llp{6.4cm}@{}}
\toprule
\textbf{Specification} & \textbf{Invariants} & \textbf{Properties} \\
\midrule
\texttt{sotfs\_graph.tla} & 8 &
  \textsc{LinkCount}, \textsc{UniqueNames},
  \textsc{DirSelfRef}, \textsc{NoDangling},
  \textsc{BlockRefcount}, \textsc{NoDirCycles},
  \textsc{ReachableHaveLinks}, \textsc{TypeInvariant} \\
\addlinespace
\texttt{sotfs\_transactions.tla} & 6 &
  \textsc{Isolation}, \textsc{Atomicity},
  \textsc{NoStaleLocks}, \textsc{RollbackCorrect},
  \textsc{TierMapping}, \textsc{TypeInvariant} \\
\addlinespace
\texttt{sotfs\_capabilities.tla} & 6 &
  \textsc{MonotonicAttenuation}, \textsc{GrantsConsistency},
  \textsc{DelegationForest}, \textsc{NoForgery},
  \textsc{TransitiveRevocation}, \textsc{TypeInvariant} \\
\addlinespace
\texttt{sotfs\_crash.tla} & 4 &
  \textsc{RecoveryConsistency},
  \textsc{NoPartialApplication},
  \textsc{RollbackToSnapshot}, \textsc{TypeInvariant} \\
\midrule
\textbf{Total} & \textbf{24} & \\
\bottomrule
\end{tabular}
\end{table}

All 24 invariants pass TLC model checking with the default state space
configuration (3 nodes, 4 edges, 2 transactions, 3 capabilities). We
note that TLA$^+$ model checking provides bounded verification---it
explores all reachable states within the specified bounds but does not
constitute a proof for unbounded instances.

\subsection{Layer 2: Crash Refinement}
\label{sec:verif-crash}

The crash refinement layer bridges TLA$^+$ specifications and the
Rust implementation. We define a \emph{crash refinement relation}
$R_{\mathrm{crash}}$ between the TLA$^+$ transaction specification
and the WAL-based implementation:

\begin{definition}[Crash Refinement]
\label{def:crash-refinement}
An implementation $I$ \emph{crash-refines} a specification $S$ if for
every execution trace $\sigma$ of $I$ (including traces interrupted by
crashes at any point), there exists a trace $\sigma'$ of $S$ such that
the observable state after recovery in $\sigma$ equals the observable
state after the corresponding recovery step in $\sigma'$.
\end{definition}

The \sotfs{} WAL ensures that every committed transaction's effects
are durable and every uncommitted transaction's effects are rolled
back upon recovery. The crash specification
(\texttt{sotfs\_crash.tla}) models arbitrary crash points during
commit sequences and verifies that recovery always reaches a
consistent state.

\subsection{Layer 3: Rust Typestate Enforcement}
\label{sec:verif-typestate}

At the implementation level, Rust's ownership and type system enforces
structural invariants statically. We use the typestate
pattern~\cite{kadekodi2024} to encode transaction lifecycle states:

\begin{lstlisting}[caption={Typestate encoding for transactions.},label=lst:typestate]
pub struct Transaction<S: TxState> {
    id: TxId,
    wal: WalHandle,
    _state: PhantomData<S>,
}

pub struct Idle;
pub struct Active;
pub struct Prepared;
pub struct Committed;

impl Transaction<Idle> {
    pub fn begin(self) -> Transaction<Active> { /* ... */ }
}
impl Transaction<Active> {
    pub fn prepare(self) -> Transaction<Prepared> { /* ... */ }
    pub fn rollback(self) -> Transaction<Idle> { /* ... */ }
}
impl Transaction<Prepared> {
    pub fn commit(self) -> Transaction<Committed> { /* ... */ }
    pub fn rollback(self) -> Transaction<Idle> { /* ... */ }
}
\end{lstlisting}

The move semantics of Rust ensure that a \texttt{Transaction<Active>}
cannot be used after calling \texttt{prepare()}---the old value is
consumed, and only the returned \texttt{Transaction<Prepared>} can
proceed. This statically prevents protocol violations such as
committing without preparing, or using a transaction after rollback.

\subsection{Layer 4: Coq Extraction (Planned)}
\label{sec:verif-coq}

The fourth layer, currently on the roadmap, targets machine-checked
proofs in Coq. The plan follows the FSCQ~\cite{chen2015}
methodology: define the graph operations in Gallina, prove the DPO
invariants (Section~\ref{sec:sotfs-invariants}) as Coq lemmas, and
extract verified Rust code via the \texttt{coq-of-rust} toolchain.
We estimate this effort at 6--12 person-months for the core graph
operations.

The planned Coq formalization targets three key theorems:

\begin{enumerate}
\item \textbf{Invariant preservation:} For each DPO rule $p$ and
  each invariant $I$ from Section~\ref{sec:sotfs-invariants}, prove
  that if $I(G)$ holds before applying $p$, then $I(H)$ holds after.
  This lifts the TLA$^+$ bounded verification to an unbounded proof.

\item \textbf{Crash refinement:} Prove that the WAL-based
  implementation crash-refines the TLA$^+$ transaction specification
  (Definition~\ref{def:crash-refinement}). This follows the
  Yggdrasil~\cite{sigurbjarnarson2016} approach of encoding the
  refinement relation as a simulation proof.

\item \textbf{Capability soundness:} Prove that no execution trace
  of the capability system can produce a capability with rights
  exceeding those of its root ancestor (Theorem~\ref{thm:monotonic}
  and Theorem~\ref{thm:revocation} in full generality).
\end{enumerate}

\subsection{Verification Summary}
\label{sec:verif-summary}

Table~\ref{tab:verif-layers} summarizes the four verification layers
and their current status.

\begin{table}[H]
\centering
\caption{Four-layer verification chain: scope and status.}
\label{tab:verif-layers}
\small
\begin{tabular}{@{}llll@{}}
\toprule
\textbf{Layer} & \textbf{Tool} & \textbf{Scope} & \textbf{Status} \\
\midrule
1. Model checking & TLA$^+$/TLC & Protocol correctness (24 inv.) & Complete \\
2. Crash refinement & TLA$^+$/SMT & WAL $\sqsubseteq$ Spec & Complete \\
3. Typestate & Rust compiler & State machine discipline & Complete \\
4. Machine proofs & Coq & End-to-end correctness & Planned \\
\bottomrule
\end{tabular}
\end{table}

The gap between layers 2--3 and layer~4 is significant: TLA$^+$
provides bounded verification (all states within the model's
parameter bounds), Rust typestate provides compile-time enforcement
of state machine discipline, but neither constitutes a proof that
the implementation matches the specification for all possible inputs
and states. Layer~4 is designed to close this gap.

% =============================================================================
\section{Implementation}
\label{sec:impl}
% =============================================================================

\subsection{Codebase Structure}
\label{sec:impl-codebase}

\sotbsd{} is implemented entirely in Rust (nightly toolchain,
\texttt{x86\_64-unknown-none} target) with minimal C glue for
bootloader integration. The codebase comprises approximately 108{,}000
lines of Rust across 69 crates, organized as follows:

\begin{itemize}[leftmargin=1.4em]
\item \textbf{\texttt{kernel/}} ($\sim$10{,}000 LOC): The five-primitive
  microkernel, including GDT/IDT, frame allocator, slab allocator,
  capability tables, SMP support (up to 4+ cores), LAPIC timer,
  SYSCALL/SYSRET fast path, and the scheduler.

\item \textbf{\texttt{libs/sotos-common/}}: Shared kernel--userspace
  ABI, syscall definitions, Linux ABI constants (198 syscall numbers,
  \texttt{Stat}, \texttt{Termios}, \texttt{Winsize},
  \texttt{KernelSigaction} structs).

\item \textbf{\texttt{libs/sotos-objstore/}}: The ObjectStore
  implementation with WAL, graph-structured metadata, block management,
  and CoW snapshots.

\item \textbf{\texttt{services/init/}}: The init process implementing
  the LUCAS (Linux User-space Compatibility and Abstraction System)
  personality, including the ELF loader, fork/exec/waitpid, signal
  delivery, vDSO injection, and 127 Linux syscall handlers.

\item \textbf{\texttt{personality/bsd/}}: BSD personality crates
  including rump kernel integration and bhyve-style VT-x
  virtualization.

\item \textbf{\texttt{personality/common/deception/}}: Deception engine
  with anomaly detection, migration orchestrator, and four deception
  profiles.

\item \textbf{\texttt{formal/}}: TLA$^+$ specifications and TLC
  configuration files.
\end{itemize}

\subsection{Key Implementation Decisions}
\label{sec:impl-decisions}

\paragraph{No heap in kernel.}
The kernel uses only fixed-size arrays and a slab allocator for
kernel objects. This eliminates heap fragmentation and simplifies
reasoning about memory safety in the TCB.

\paragraph{Spinlocks throughout.}
All kernel data structures are protected by ticket spinlocks
(\texttt{spin::Mutex}). While this limits scalability on many-core
systems, it simplifies the concurrency model and avoids priority
inversion issues in the current prototype.

\paragraph{Copy-on-Write fork.}
The \texttt{SYS\_AS\_CLONE} syscall implements true CoW fork by
walking all user page table entries, marking parent and child
read-only, and incrementing per-frame reference counts. Write faults
trigger frame duplication in the userspace VMM.

\paragraph{vDSO injection.}
Each child process receives a forged vDSO ELF at virtual address
\texttt{0xB80000} providing \texttt{\_\_vdso\_clock\_gettime}
(RDTSC-based), signal trampolines, and fork-restore trampolines.
The vDSO is a complete ELF shared library with \texttt{PT\_DYNAMIC},
\texttt{.dynsym}, \texttt{.gnu.hash}, and \texttt{LINUX\_2.6}
version definitions.

\subsection{Filesystem Implementation}
\label{sec:impl-fs}

The \sotfs{} implementation stores the typed graph in a two-level
structure:

\begin{enumerate}
\item \textbf{Node table:} A fixed-size array of node descriptors,
  each containing the node type, attributes, and an adjacency list
  (stored as a sorted array of edge indices).

\item \textbf{Edge table:} A fixed-size array of edge descriptors,
  each containing source, target, edge type, name (for
  \texttt{contains} edges), and attributes.

\item \textbf{Block layer:} Data blocks are stored in a bitmap-managed
  region with CoW semantics for snapshot support.
\end{enumerate}

DPO rule application is implemented as atomic compare-and-swap
sequences on the node and edge tables, wrapped in WAL transactions
for crash consistency.

% =============================================================================
\section{Evaluation}
\label{sec:eval}
% =============================================================================

We evaluate \sotfs{} along five dimensions: DPO operation latency,
data I/O throughput, monitor scaling, correctness testing, and
comparison with existing filesystems.

All benchmarks were collected on a QEMU \texttt{x86\_64} virtual
machine with 512\,MiB RAM, running on an Intel Core i7-12700H host.
Timing uses RDTSC with TSC frequency calibrated via CPUID leaf 15H.

\subsection{DPO Operation Latency}
\label{sec:eval-dpo}

Table~\ref{tab:dpo-latency} reports the latency of DPO rewriting
operations at two scales. Per-operation costs are dominated by WAL
logging and graph invariant checking.

\begin{table}[H]
\centering
\caption{DPO operation latency. Per-op is amortized over the larger batch.}
\label{tab:dpo-latency}
\small
\begin{tabular}{@{}lrrr@{}}
\toprule
\textbf{Operation} & \textbf{100 ops} & \textbf{500 ops} & \textbf{Per-op (amortized)} \\
\midrule
\texttt{create\_file} & 84\,$\mu$s  & 1.8\,ms  & $\sim$9.6\,$\mu$s \\
\texttt{mkdir} (chain) & 36\,$\mu$s (50) & 622\,$\mu$s & $\sim$1.2\,$\mu$s \\
\texttt{unlink}        & 104\,$\mu$s & 1.9\,ms  & $\sim$3.8\,$\mu$s \\
\texttt{rename}        & 255\,$\mu$s & 5.8\,ms  & $\sim$11.6\,$\mu$s \\
\texttt{hard\_link}    & 2.7\,$\mu$s (10) & 1.9\,ms (500) & $\sim$3.9\,$\mu$s \\
\bottomrule
\end{tabular}
\end{table}

The \texttt{rename} operation is the most expensive because it
performs two graph modifications (edge deletion and edge creation)
atomically within a single DPO rule application, requiring both the
dangling condition check on the source directory and the unique-name
check on the destination directory. The \texttt{mkdir} chain benchmark
creates directories $d_1/d_2/\ldots/d_n$ in sequence, showing linear
scaling with a low constant (1.2\,$\mu$s per directory) because each
operation touches only the parent--child edge. The \texttt{hard\_link}
operation at 10 links takes only 2.7\,$\mu$s because it adds a single
edge without modifying existing structure; at 500 links the amortized
cost rises to 3.9\,$\mu$s due to treewidth maintenance overhead as the
hub node's bag grows.

\paragraph{Scaling analysis.}
The roughly linear scaling from 100 to 500 operations (e.g.,
\texttt{create\_file}: $84\,\mu\text{s} \to 1.8\,\text{ms}$, a
factor of 21$\times$ for 5$\times$ more operations) suggests a
mild super-linear component attributable to WAL log scanning
during invariant verification. We expect this to improve with a
hash-indexed WAL.

\subsection{Data I/O Performance}
\label{sec:eval-io}

Table~\ref{tab:io} reports raw read and write latencies for data
blocks of varying sizes.

\begin{table}[H]
\centering
\caption{Data I/O latency by block size.}
\label{tab:io}
\small
\begin{tabular}{@{}lrr@{}}
\toprule
\textbf{Size} & \textbf{Write} & \textbf{Read} \\
\midrule
64\,B  & 929\,ns & 54\,ns \\
1\,KB  & 782\,ns & 60\,ns \\
4\,KB  & 805\,ns & 71\,ns \\
64\,KB & 3.5\,$\mu$s & --- \\
\bottomrule
\end{tabular}
\end{table}

Read latencies are dominated by memory copy cost and scale linearly
with size. Write latencies include WAL entry creation. The anomalous
result for 64\,B writes (929\,ns $>$ 1\,KB writes at 782\,ns)
reflects WAL entry overhead being amortized over larger payloads.

\subsection{Monitor Scaling}
\label{sec:eval-monitors}

Table~\ref{tab:monitors} reports the computation time for the
treewidth estimator and Ollivier--Ricci curvature monitor as the
graph size increases. We use star topologies (one root, $n-1$ leaves)
to stress the curvature computation and chain topologies
($v_1 \to v_2 \to \cdots \to v_n$) as a lower bound.

\begin{table}[H]
\centering
\caption{Monitor computation time by graph size and topology.}
\label{tab:monitors}
\small
\begin{tabular}{@{}lrrrr@{}}
\toprule
\textbf{Monitor (topology)} & \textbf{50 nodes} & \textbf{100 nodes} & \textbf{200 nodes} & \textbf{500 nodes} \\
\midrule
Treewidth (star)    & 46\,$\mu$s  & 192\,$\mu$s  & 917\,$\mu$s  & 6.4\,ms \\
Curvature (star)    & 55\,$\mu$s  & 138\,$\mu$s  & 450\,$\mu$s  & 2.7\,ms \\
Curvature (chain)   & 66\,$\mu$s  & 159\,$\mu$s  & 302\,$\mu$s  & --- \\
\bottomrule
\end{tabular}
\end{table}

The treewidth estimator scales super-linearly on star topologies due
to the $O(n \cdot k)$ cost of maintaining tree decomposition bags
where $k$ is the current width. Curvature computation scales
better because the Wasserstein distance for each edge depends only on
local neighborhood structure (1-hop neighborhoods), giving
$O(|E| \cdot \Delta^2)$ total cost where $\Delta$ is the maximum
degree. For chains ($\Delta = 2$), curvature is nearly linear.

For a typical filesystem with 10{,}000 files organized in a balanced
directory hierarchy (depth 4--5, branching factor 20--30), the
treewidth is at most 6 (no hard links) and curvature computation
completes in under 50\,ms, well within the monitoring epoch of
1 second used in our anomaly detection pipeline.

\paragraph{Asymptotic analysis.}
Let $n = |V|$, $m = |E|$, and $\Delta = \max_{v \in V} \deg(v)$.

\begin{itemize}[leftmargin=1.4em]
\item \textbf{Treewidth estimation:} We use a greedy heuristic
  (minimum-degree elimination) that runs in $O(n \cdot \Delta)$ time.
  For star topologies, $\Delta = n - 1$, giving $O(n^2)$ and
  explaining the super-linear scaling in Table~\ref{tab:monitors}
  ($46\,\mu\text{s} \to 192\,\mu\text{s} \to 917\,\mu\text{s}
  \to 6.4\,\text{ms}$ for $50 \to 100 \to 200 \to 500$ nodes).

\item \textbf{Curvature computation:} For each edge $(u, v)$, we
  solve a linear program of size $O(\deg(u) \cdot \deg(v))$ to
  compute the Wasserstein distance. Total cost is
  $O(m \cdot \Delta^2)$. For chains ($\Delta = 2$), this is $O(n)$;
  for stars ($\Delta = n - 1$), this is $O(n^2)$, matching the
  observed scaling.
\end{itemize}

\paragraph{Incremental update potential.}
When a single DPO rule modifies $O(1)$ edges, only the curvature
values of edges in the 2-hop neighborhood of the modification need
recomputation. This reduces incremental update cost to
$O(\Delta^3)$ per operation---constant for bounded-degree graphs.
We have not yet implemented incremental updates but estimate they
would reduce monitoring overhead by 2--3 orders of magnitude for
large graphs.

\subsection{Correctness and Robustness}
\label{sec:eval-correctness}

We evaluate correctness through four complementary methodologies:

\paragraph{Unit and integration tests.}
76 tests covering all DPO rules, graph invariants, transaction
semantics, capability operations, and curvature computation.
All 76 pass with zero failures.

\paragraph{Property-based testing (proptest).}
7 property tests, each executing 100 random operation sequences
(create, delete, rename, link, mkdir, rmdir, chmod in random order
on random targets). Two bugs were discovered and fixed:
\texttt{rename(".")} was incorrectly permitted (violating the
directory self-reference invariant), and \texttt{link("..")} created
an illegal edge violating type constraints.

\paragraph{Fuzz testing (cargo-fuzz).}
3 fuzz targets (graph operations, transaction sequences, capability
chains) each ran for 100{,}000 iterations, totaling 300{,}000 fuzz
corpus runs. Zero crashes were observed, and no invariant violations
were detected.

\paragraph{Adversarial testing.}
13 attack scenarios were executed against the full \sotfs{} stack:

\begin{table}[H]
\centering
\caption{Adversarial test scenarios and detection results.}
\label{tab:adversarial}
\small
\begin{tabular}{@{}lll@{}}
\toprule
\textbf{Scenario} & \textbf{Detection Method} & \textbf{Result} \\
\midrule
Ransomware (mass encrypt)        & Curvature drop ($\Delta\OR < -0.15$) & Detected \\
Ransomware (rename-in-place)     & Curvature + rename rate               & Detected \\
Hard-link bomb                   & Treewidth spike ($\tw > 10$)          & Detected \\
Deep directory chain             & Curvature (chain pattern)              & Detected \\
Mass deletion                    & Curvature + edge count drop            & Detected \\
Exfiltration staging             & Fan-out anomaly                        & Detected \\
Permission escalation (chmod)    & Capability graph inconsistency         & Detected \\
Backdoor file creation           & Honeypot trigger                       & Detected \\
TOCTOU rename attack             & DPO atomicity (no race window)         & Prevented \\
Symlink traversal                & Type constraint violation              & Prevented \\
Cross-directory hard link escape & Capability scope check                 & Prevented \\
Inode exhaustion DoS             & Resource limit enforcement             & Prevented \\
Transaction starvation           & Lock timeout + \textsc{NoStaleLocks}   & Prevented \\
\bottomrule
\end{tabular}
\end{table}

All 13 scenarios were either detected (curvature/treewidth alert
raised) or prevented (operation rejected by invariant enforcement).

\paragraph{Kernel QA.}
Three kernel-level stress tests verify the foundation:

\begin{itemize}[leftmargin=1.4em]
\item \textbf{Capability escalation:} 10 attempts to forge or
  amplify capabilities, all rejected (10/10 PASS).
\item \textbf{IPC storm:} 4 threads sending 1{,}000 messages each
  concurrently, all delivered correctly (4/4 PASS).
\item \textbf{SMP atomics stress:} 4{,}000 atomic
  increment/decrement operations across 4 cores, final count
  correct (4{,}000/4{,}000 PASS).
\end{itemize}

A kernel bug was discovered during stress testing: the endpoint pool
was limited to 64 slots per core, causing exhaustion under IPC storm.
This was fixed by introducing \texttt{SYS\_ENDPOINT\_CLOSE} and
increasing \texttt{MAX\_EP} to 256.

\subsection{Comparison with Existing Filesystems}
\label{sec:eval-comparison}

Table~\ref{tab:comparison} compares \sotfs{} with existing filesystems
across the dimensions that motivate this work. We emphasize that
\sotfs{} is a research prototype and does not yet match the raw
throughput of production filesystems; the comparison highlights
\emph{architectural} capabilities.

\begin{table}[H]
\centering
\caption{Feature comparison with existing filesystems. A dash (---)
  indicates the feature is absent.}
\label{tab:comparison}
\small
\begin{tabular}{@{}lcccccc@{}}
\toprule
\textbf{Feature} & \textbf{ext4} & \textbf{ZFS} & \textbf{Btrfs} & \textbf{HAMMER2} & \textbf{SquirrelFS} & \textbf{sotFS} \\
\midrule
Metadata model & tables & Merkle & B-trees & radix & Rust types & typed graph \\
Formal ops     & --- & --- & --- & --- & typestate & DPO rules \\
Treewidth bound & --- & --- & --- & --- & --- & $\leq k$ \\
MSO queries    & --- & --- & --- & --- & --- & $O(n)$ \\
Anomaly detect & --- & --- & --- & --- & --- & Ollivier--Ricci \\
Cap integration & DAC & DAC & DAC & DAC & --- & graph edges \\
Deception      & --- & --- & --- & --- & --- & subgraph views \\
Verification   & tested & tested & tested & tested & typestate & TLA$^+$$\to$Coq \\
\bottomrule
\end{tabular}
\end{table}

\sotfs{} is, to our knowledge, the only filesystem with a formal
algebraic operational semantics (DPO rules), bounded treewidth
guarantees enabling efficient MSO queries, topology-based anomaly
detection, and integrated deception capabilities.

\paragraph{Comparison methodology.}
The features compared in Table~\ref{tab:comparison} represent
architectural capabilities, not performance characteristics.
Production filesystems (ext4, ZFS, Btrfs) have decades of
engineering investment and vastly superior raw throughput. Our
claim is not that \sotfs{} is faster, but that its graph-structured
metadata model enables capabilities that are structurally impossible
in flat-metadata systems, regardless of engineering effort.
Specifically:

\begin{itemize}[leftmargin=1.4em]
\item \textbf{MSO queries} require bounded treewidth, which in turn
  requires explicit graph structure. Ext4's inode tables have no
  natural tree decomposition with bounded width.

\item \textbf{Ollivier--Ricci curvature} is defined on graphs. While
  one could construct a graph from ext4 metadata post-hoc, this would
  be an offline analysis tool rather than an integrated runtime
  monitor.

\item \textbf{Deception via subgraph projection} requires graph
  homomorphisms and functorial consistency, which presuppose a
  graph-structured data model.
\end{itemize}

SquirrelFS~\cite{kadekodi2024} is the closest related work in
terms of leveraging Rust's type system for filesystem correctness.
However, SquirrelFS focuses on crash consistency via typestate
patterns, while \sotfs{} provides a broader algebraic framework
that additionally enables structural queries and anomaly detection.
The approaches are complementary: \sotfs{}'s typestate layer
(Section~\ref{sec:verif-typestate}) draws directly on SquirrelFS's
technique.

% =============================================================================
\section{Discussion}
\label{sec:discussion}
% =============================================================================

\subsection{Limitations}
\label{sec:discussion-limits}

\paragraph{Scalability.}
The current implementation uses spinlocks and fixed-size arrays,
limiting scalability to small graph sizes (tested up to 500 nodes
in benchmarks). A production deployment would require replacing
spinlocks with RCU-style read-side critical sections and supporting
dynamic array growth.

\paragraph{Treewidth with hard links.}
The $\mathrm{LINK\_MAX} + O(1)$ treewidth bound with hard links is
technically vacuous for systems with $\mathrm{LINK\_MAX} = 65{,}535$:
treewidth that large provides no practical benefit for Courcelle's
theorem, since the constant factor grows exponentially with treewidth.
In practice, we expect most filesystem graphs to have treewidth under
20 (few files have more than a handful of hard links), but we have not
yet characterized real-world treewidth distributions.

\paragraph{Coq verification gap.}
Layer 4 of our verification chain (Coq extraction) is planned but not
yet implemented. The current guarantee is bounded model checking
(TLA$^+$) plus implementation-level typestate enforcement, which does
not constitute a machine-checked end-to-end proof.

\paragraph{Curvature computation cost.}
While curvature computation is efficient for moderate-sized graphs
($<$50\,ms for 10{,}000 nodes), it becomes expensive for very large
filesystems with millions of files. Incremental curvature updates
(recomputing only affected edges after each operation) would reduce
this cost but are not yet implemented.

\paragraph{Performance gap.}
\sotfs{} runs as a userspace service communicating with the kernel
via IPC, adding per-operation overhead compared to in-kernel
filesystems. The 9.6\,$\mu$s per \texttt{create\_file} is acceptable
for a research prototype but is roughly 10$\times$ slower than ext4's
metadata path on the same hardware.

\subsection{Applicability Beyond Filesystems}
\label{sec:discussion-applicability}

The graph-structured approach with DPO rewriting and curvature
monitoring is not inherently filesystem-specific. Any system whose
state can be modeled as a typed graph---network topologies, access
control graphs, service dependency graphs, software module
graphs---could benefit from the same framework. The \sotbsd{} SOT
layer already represents all kernel objects as typed secure objects
with capability edges, suggesting that the monitoring infrastructure
could be extended to cover the entire system state, not just
filesystem metadata.

In particular, the deception functor framework
(Section~\ref{sec:deception-functor}) generalizes to any
capability-mediated system: the deceptive view is simply a subgraph
projection filtered by the observer's capability set. This could
enable deception at the network level (fake service topologies),
the process level (decoy processes), or the memory level (honeypot
pages), all within a unified categorical framework.

\subsection{Threat Model Assumptions}
\label{sec:discussion-threats}

Our evaluation assumes an attacker who has gained code execution
within a sandboxed process but does not have access to the kernel
TCB. The attacker can invoke any syscall, open any file reachable
through their capabilities, and perform arbitrary computation. The
deception engine assumes the attacker cannot distinguish the true
filesystem from the deceptive projection without performing
operations that trigger curvature alerts.

This threat model does not cover: (i)~hardware attacks (DMA, cold
boot), (ii)~side-channel attacks (cache timing, speculative
execution), or (iii)~compromised kernel code. The first two are
orthogonal to our contributions; the third requires the end-to-end
Coq verification that is currently planned but not yet implemented.

\subsection{Future Work}
\label{sec:discussion-future}

Five directions are immediate priorities:

\begin{enumerate}
\item \textbf{Coq formalization:} Encode the DPO rules and graph
  invariants in Coq, prove preservation theorems, and extract
  verified Rust code. We estimate 6--12 person-months for the core
  graph operations, following the FSCQ methodology.

\item \textbf{Incremental monitors:} Maintain treewidth
  decompositions and curvature values incrementally rather than
  recomputing from scratch, reducing monitor overhead to $O(\Delta^3)$
  amortized per operation for curvature and $O(\Delta)$ for treewidth.

\item \textbf{Real-world treewidth characterization:} Instrument
  ext4/XFS/Btrfs installations to measure treewidth distributions
  of real filesystem graphs, validating our assumption that practical
  treewidth is small. We plan to analyze filesystem snapshots from
  production servers with $10^6$--$10^8$ files.

\item \textbf{Distributed \sotfs{}:} Extend the DPO framework to
  distributed graphs spanning multiple nodes, using the categorical
  framework to define distributed transactions with crash-refinement
  guarantees.

\item \textbf{GPU-accelerated curvature:} The Wasserstein distance
  computation for each edge is an independent linear program, making
  it embarrassingly parallel. A GPU implementation could enable
  real-time curvature monitoring for graphs with millions of nodes.
\end{enumerate}

% =============================================================================
\section{Conclusion}
\label{sec:conclusion}
% =============================================================================

We have presented \sotbsd{}, a capability-secure microkernel operating
system, and \sotfs{}, a graph-structured filesystem with algebraic
operational semantics. By making the graph structure of filesystem
metadata explicit and equipping POSIX operations with DPO rewriting
rules, we simultaneously achieve three goals: linear-time MSO queries
via bounded treewidth and Courcelle's theorem, topology-aware anomaly
detection through Ollivier--Ricci curvature, and capability-gated
deceptive projections.

Our evaluation demonstrates that the algebraic approach is practical:
DPO operations complete in under 12\,$\mu$s, monitors scale to 500
nodes in milliseconds, and a comprehensive testing campaign (76 tests,
300{,}000 fuzz runs, 13 adversarial scenarios) reveals no crashes or
undetected attacks. The four-layer verification chain provides
defense-in-depth against specification, protocol, and implementation
bugs, with machine-checked Coq proofs as the planned capstone.

\sotfs{} demonstrates that principled mathematical
foundations---category theory, graph rewriting, discrete differential
geometry---can be brought to bear on practical filesystem design,
yielding capabilities (MSO queries, curvature-based anomaly detection,
functorial deception) that are impossible in systems without explicit
graph structure.

More broadly, this work suggests a design principle for security-critical
systems: \emph{make implicit structure explicit}. Filesystem metadata
has always been a graph; by representing it as one, we unlock an entire
mathematical toolkit---tree decompositions, Ricci curvature, graph
homomorphisms, functors---that was previously inaccessible. The same
principle applies to network topologies, access control policies,
and software architectures: wherever structure is implicit, making it
explicit enables formal reasoning, efficient querying, and principled
security monitoring.

The \sotbsd{} microkernel and \sotfs{} filesystem are open-source and
available for reproduction of all results reported in this paper.

% =============================================================================
% REFERENCES
% =============================================================================
\newpage
\begin{thebibliography}{30}

\bibitem{ehrig2006}
H.~Ehrig, K.~Ehrig, U.~Prange, and G.~Taentzer,
\emph{Fundamentals of Algebraic Graph Transformation},
Springer, 2006.

\bibitem{courcelle1990}
B.~Courcelle,
``The monadic second-order logic of graphs. I. Recognizable sets of
finite graphs,''
\emph{Information and Computation}, vol.~85, no.~1, pp.~12--75, 1990.

\bibitem{ollivier2009}
Y.~Ollivier,
``Ricci curvature of Markov chains on metric spaces,''
\emph{Journal of Functional Analysis}, vol.~256, no.~3,
pp.~810--864, 2009.

\bibitem{lin2011}
Y.~Lin, L.~Lu, and S.-T.~Yau,
``Ricci curvature of graphs,''
\emph{Tohoku Mathematical Journal}, vol.~63, no.~4,
pp.~605--627, 2011.

\bibitem{chen2015}
H.~Chen, D.~Ziegler, T.~Chajed, A.~Chlipala, M.~F.~Kaashoek,
and N.~Zeldovich,
``Using Crash Hoare Logic for certifying the FSCQ file system,''
in \emph{Proc.\ 25th ACM SOSP}, 2015, pp.~18--37.

\bibitem{sigurbjarnarson2016}
H.~Sigurbjarnarson, J.~Bornholt, E.~Torlak, and X.~Wang,
``Push-button verification of file systems via crash refinement,''
in \emph{Proc.\ 12th USENIX OSDI}, 2016, pp.~1--16.

\bibitem{kadekodi2024}
S.~Kadekodi, T.~Pillai, R.~Alagappan, and A.~C.~Arpaci-Dusseau,
``SquirrelFS: Using the Rust compiler to check file-system
crash consistency,''
in \emph{Proc.\ USENIX OSDI}, 2024.

\bibitem{watson2010}
R.~N.~M.~Watson, J.~Anderson, B.~Laurie, and K.~Kennaway,
``Capsicum: Practical capabilities for UNIX,''
in \emph{Proc.\ 19th USENIX Security}, 2010, pp.~29--46.

\bibitem{klein2009}
G.~Klein, K.~Elphinstone, G.~Heiser, J.~Andronick, D.~Cock,
P.~Derrin, D.~Elkaduwe, K.~Engelhardt, R.~Kolanski, M.~Norrish,
T.~Sewell, H.~Tuch, and S.~Winwood,
``seL4: Formal verification of an OS kernel,''
in \emph{Proc.\ 22nd ACM SOSP}, 2009, pp.~207--220.

\bibitem{chajed2019}
T.~Chajed, J.~Tassarotti, M.~F.~Kaashoek, and N.~Zeldovich,
``Verifying concurrent, crash-safe systems with Perennial,''
in \emph{Proc.\ 27th ACM SOSP}, 2019, pp.~243--258.

\bibitem{robertson1986}
N.~Robertson and P.~D.~Seymour,
``Graph minors. II. Algorithmic aspects of tree-width,''
\emph{Journal of Algorithms}, vol.~7, no.~3, pp.~309--322, 1986.

\bibitem{bodlaender1998}
H.~L.~Bodlaender,
``A partial $k$-arboretum of graphs with bounded treewidth,''
\emph{Theoretical Computer Science}, vol.~209, no.~1--2,
pp.~1--45, 1998.

\bibitem{woodruff2014}
J.~Woodruff, R.~N.~M.~Watson, D.~Chisnall, S.~W.~Moore,
J.~Anderson, B.~Davis, B.~Laurie, P.~G.~Neumann, R.~Norton,
and M.~Roe,
``The CHERI capability model: Revisiting RISC in an age of risk,''
in \emph{Proc.\ 41st ACM/IEEE ISCA}, 2014, pp.~457--468.

\bibitem{feske2015}
N.~Feske,
``Genode Operating System Framework,''
\emph{Genode Labs}, 2015.
\url{https://genode.org}

\bibitem{liedtke1993}
J.~Liedtke,
``Improving IPC by kernel design,''
in \emph{Proc.\ 14th ACM SOSP}, 1993, pp.~175--188.

\bibitem{muniswamy2006}
K.-K.~Muniswamy-Reddy, D.~A.~Holland, U.~Braun, and M.~Seltzer,
``Provenance-aware storage systems,''
in \emph{Proc.\ USENIX ATC}, 2006, pp.~43--56.

\bibitem{jajodia2011}
S.~Jajodia, A.~K.~Ghosh, V.~Swarup, C.~Wang, and X.~S.~Wang,
\emph{Moving Target Defense: Creating Asymmetric Uncertainty for
Cyber Threats},
Springer, 2011.

\bibitem{provos2004}
N.~Provos,
``A virtual honeypot framework,''
in \emph{Proc.\ 13th USENIX Security}, 2004, pp.~1--14.

\bibitem{kantee2012}
A.~Kantee,
``Flexible operating system internals: The design and implementation
of the Anykernel and Rump Kernels,''
Ph.D. dissertation, Aalto University, 2012.

\bibitem{bonwick2003}
J.~Bonwick and B.~Moore,
``ZFS: The last word in file systems,''
\emph{Sun Microsystems}, 2003.

\bibitem{rodeh2013}
O.~Rodeh, J.~Bacik, and C.~Mason,
``BTRFS: The Linux B-Tree Filesystem,''
\emph{ACM Transactions on Storage}, vol.~9, no.~3, 2013.

\bibitem{jung2017}
R.~Jung, J.-H.~Jourdan, R.~Krebbers, and D.~Dreyer,
``RustBelt: Securing the foundations of the Rust programming
language,''
\emph{Proc.\ ACM on Programming Languages}, vol.~2, no.~POPL,
pp.~66:1--66:34, 2017.

\bibitem{pillai2014}
T.~S.~Pillai, V.~Chidambaram, R.~Alagappan, S.~Al-Kiswany,
A.~C.~Arpaci-Dusseau, and R.~H.~Arpaci-Dusseau,
``All file systems are not created equal: On the complexity of
crafting crash-consistent applications,''
in \emph{Proc.\ 11th USENIX OSDI}, 2014, pp.~433--448.

\bibitem{jannen2015}
W.~Jannen, J.~Yuan, Y.~Zhan, A.~Akshintala, J.~Esmet, Y.~Jiao,
A.~Mittal, P.~Pandey, P.~Reddy, L.~Walsh, M.~Bender, M.~Farach-Colton,
R.~Johnson, B.~C.~Kuszmaul, and D.~E.~Porter,
``BetrFS: A right-optimized write-optimized file system,''
in \emph{Proc.\ 13th USENIX FAST}, 2015, pp.~301--315.

\bibitem{king2003}
S.~T.~King and P.~M.~Chen,
``Backtracking intrusions,''
in \emph{Proc.\ 19th ACM SOSP}, 2003, pp.~223--236.

\bibitem{goel2005}
A.~Goel, K.~Po, K.~Farhadi, Z.~Li, and E.~de~Lara,
``The Taser intrusion recovery system,''
in \emph{Proc.\ 20th ACM SOSP}, 2005, pp.~163--176.

\bibitem{dennis1966}
J.~B.~Dennis and E.~C.~Van~Horn,
``Programming semantics for multiprogrammed computations,''
\emph{Communications of the ACM}, vol.~9, no.~3, pp.~143--155, 1966.

\bibitem{dillon2015}
M.~Dillon,
``HAMMER2: The next-generation filesystem for DragonFly BSD,''
\emph{DragonFly BSD Project}, 2015.

\bibitem{illumos2025}
{illumos Project},
``Service Management Facility (SMF),''
\url{https://illumos.org/man/7/smf}, accessed April 2026.

\end{thebibliography}

\end{document}
